\documentclass[12pt]{article}
\usepackage{color,soul}
% Import preambles and macros for notes
\input{../../../../preambles/notes-preamble.tex}
\input{../../../../preambles/notes-macros.tex}
\input{../../../../preambles/theorem-system-section.tex}

\title{MATH 420 Lecture 14}
\author{Deepak Jassal}
\date{March 4\textsuperscript{th}, 2026}

\begin{document}
\maketitle
\stepcounter{section}
\section*{Modules}
They are generalizations of vector spaces.
\defn{Modules}{
    Let $R$ be a ring. An Abelian group $(M,\oplus)$ is said to be a left-$R$ module with respect to an operation
    \[
        \dot:R\times M\to M
    \]
    if given $r\in R, x,y\in M$
    \begin{enumerate}
        \item $r\dot(x\oplus y)=r\dot x\oplus r\dot y$,$(r+ y)\dot x=r\dot x\oplus y\dot x$;
        \item $r\dot(x\dot y)=(r\dot x)\dot y$;
        \item $1\dot x=x$.
    \end{enumerate}
}
All Abelian groups are $\Z-$modules.
\ex{
    $\Z^3$ is a $\Z-$module.\\
}
Any subgroup of $(\Z^n,+)$ is a $\Z-$module, such a subgroup is denoted by $\Lambda$, say (lattice). The rank of $\Lambda$ is the max number of linearly independant vectors in $\Lambda$.\\
$\rank(\Z^3)=3$.\\
Q: Given $\Lambda\subset\Z^n$, compute a basis of a $\Lambda$. i.e., $B\subset\Lambda$ such that $\spn_{\Z}(B)=\Lambda$, $|B|=\rank(\Lambda)$.\\
Shortest basis problem for $\Lambda\subset\Z^n$ is NP-hard.\\
If $\Lambda=\spn_Z(B),B=\{\vec{b_1},\dots,\vec{b_n}\}\subset\Z^n$, then a theorem of Minkowski assers that there exists a possitive number $f(a)$ such that the shortest vector $\vec{v}\in\Lambda$ has length at most $f(n)(\det\Lambda)^{\frac{1}{n}}$, $\det\Lambda=|\det(B)|$.
\defn{Rank of a Ring}{
    A commutative ring is said to have an \underline{order} rank of $r$ if $(R,+)$ is isomorphic to $\Z^r$ as a $\Z-$module, $R$ is unital.
}
Q: Given an order $R$ of a finite rank, decide if $x^2+1=0$ has a solution in $R$.\\
(Lenstra) This can be decided with an algorithm in P. But to decide $x^2+x+1=0$ is NP hard. 
\ex{
    Elliptic Cruves
    \[
        E:y^2=x^3+Ax+B,\quad A,B\in\Z,\quad 4A^3+27B^2\neq0.
    \]
    \[
        E(\Q)=\{(x,y)\in\Q^2:y^2=x^3+Ax+B\}.
    \]
    Mordell (1922): $E(\Q)\cong G\oplus\Z^r$, for some $r\geq0$, $G$ a finite Abelian group, as groups. $G$ is the torsion part and $\Z^r$ is the free part. The $r$ in Mordell's theorem is the (algebraic) rank of $E$. 
}
Tate module.\\
Fix a prime $p$ (usually $p\nmid16(4A^3-27B^2)$). Over $\C$, $E(\C)$ has a $p-$torsion point. $E(\C)$ is a group, ``torsion'' in a group means elements of finite order.\\
Let $E[P]$ denote the set of $p-$torsion points in $E(\C)$.\\
We can play the same game for $p^n,$ $n\geq1$. Get $E[p^n]$. Turns out $E[P^n]$ are \underline{algebraic} for all $p^n$.
\section*{Algebras}
\ex{
    Set of polynomials with real coefficients is an $\R-$algebra. $\R[x]$ clearly has the same structure of a $\R-$vector space is also a ring. 
}
\defn{$A$-Algebra}{
    Let $A$ be a ring. An $A$-algebra is a ring $B$ together with a homomorphism $i_B:A\to B$.
}
\defn{}{
    Let $B,C$ be $A-$algebras. A homomorphism $\phi:B\to C$ is a map such that $\phi(i_B(a))=i(a)$, for all $a\in A$, where $\phi$ is a ring homomorphism.
}
\end{document}